% Cote 119: offprints, reports, courses, CV and letters (typed and printed part): archive cote 119, site page 221 (PDF page 222).
% typescript. Draft compared: none.
% AI-transcribed by Claude Opus 5.5 (claude-opus-5-5, adaptive thinking, effort high), Message Batches API, batch r119_full (msgbatch_01DwZNm5Y7YsKBWyAN4sFs8w), 2026-11-13; not yet checked by eye.
\scanpage{119/221}
\mspage{. 8 .}

b) \uline{Théorèmes de dualité locaux et globaux pour les}
\uline{faisceaux cohérents}, développés notamment par
J.L. Verdier [$*$] et J.P. Ramis et G. Ruget [$*$], inspirés
par la théorie que j'avais développée dans le cas des
schémas, exposée dans R. Hartshorne [$*$].

c) \uline{Formulations de théorèmes du type de Riemann-Roch}
pour des variétés analytiques compactes ou des morphismes
propres de telles variétés, cf. [16, SGA 6, Exp. 0].
Les problèmes essentiels restent toujours ouverts.

d) \uline{Théorèmes de De Rham analytiques complexes} [16 bis],
\uline{cohomologie cristalline complexe}. Certains des résultats
et des idées que j'avais développés à ce sujet ont été
utilisés dans des développements théoriques divers, comme
la théorie de Hodge généralisée de P. Deligne [$*$].

e) \uline{Espaces rigide-analytiques}. M'inspirant de l'exemple
de la ``courbe elliptique de Tate'', et des besoins de
la ``géométrie formelle'' sur un anneau de valuation
discrète complet, j'étais parvenu à une formulation
partielle de la notion de variété rigide-analytique
sur un corps valué complet, qui a joué son rôle dans
la première étude systématique de cette notion par
J. Tate [$*$]. Par ailleurs, les ``cristaux'' que j'introduis
sur les variétés algébriques sur un corps de caractéris-%
tique $> 0$ peuvent s'interpréter parfois en termes de
fibrés vectoriels à connexion intégrable sur certains
types d'espaces rigide-analytiques sur des corps de
caractéristique nulle ; ceci fait pressentir l'existence
de relations profondes entre cohomologie cristalline
en car. $> 0$ , et cohomologie de systèmes locaux sur
des variétés rigide-analytiques en car. nulle.

\uline{Références} : J.P. Serre, H. Grauert, H. Cartan, P. Deligne,
A. Douady, B. Malgrange, K. Knorr, R. Kiehl,
J. Tate.

\margin{221}
